% paper/tables/metrics.tex -- generated by analyze/paper/tables_paper.py; sources in paper/tables/README.md
\begin{tabular}{@{}p{0.17\linewidth}p{0.33\linewidth}p{0.15\linewidth}p{0.26\linewidth}@{}}
\toprule
指標 & 何を測るか & 評価データ／点 & 位置づけ \\
\midrule
\multicolumn{4}{@{}l}{\textbf{自分側の見込み}} \\
自己位置の追跡 & 運動が見えるステップで $h^1$ から A-1 の位置を線形回帰で読む決定係数（0.80 以上を前提とする） & 共通／MSE & 事前登録した前提 \\
価値と行動・色の結びつき（3指標） & 運動が隠れたステップで予測が赤の方へ一歩動いた視界に寄るか（定義は補足資料 S13） & 共通／MSE & 事前登録した主 \\
自己側の積み上がり $g$ & 運動と視覚を止めた15ステップで予測が赤の方へだけ進み続ける傾きの差 & 共通／MSE & 事前登録した主 \\
快・不快への依存 $d_P$ & 快・不快を保ったときの $g$ と，符号を反転して保ったときの $g$ の差の半分 & 共通／MSE & 事前登録（符号無作為条件との比較）．好み切替条件では事前登録外 \\
\midrule
\multicolumn{4}{@{}l}{\textbf{他者側への転用}} \\
他者側の符号応答 G5b & process-2 の1つの範囲に $\pm a$ を入れたときの予測の中の A-2 のずれの差 & 共通（A-2 静止）／L1 & 事前登録した主 \\
符号に従う分 & 本当の快・不快を保った傾きと符号を反転して保った傾きの差の半分 & 緑好み，好み混合／L1 & 定義は事後（組1〜4）．組5・6 の経験の比較と符号無作為条件との比較は事前登録 \\
他者側の積み上がり G5a & 本当の快・不快を保ったときの予測の中の A-2 の傾き（戻る動きが混ざる） & 緑好み，好み混合／L1 & 事前登録した関門．結果を見た後に副へ \\
真値条件の $\Delta$ & process-2 に本当の快・不快を入れたときと零を入れたときの他者側の特徴予測の誤差の差 & 緑好み，好み混合／L1 の最後 & 事前登録外（悪化は予想外）．規則の確率の比較と符号無作為の組1〜3 は事前登録 \\
\midrule
\multicolumn{4}{@{}l}{\textbf{他者用モジュールの見分け}} \\
入れ替え応答 T2 & $t=50$ で赤と緑を入れ替えたとき $Q2$ が新しく緑が見える範囲へ移るか & 緑好み／段階II & 事前登録した主 \\
好み差応答 T5 & 位置をそろえた緑好き・赤好きの A-2 で $Q2$ が好みに応じて違うか & 緑好み，赤好み／段階II & 事前登録した主 \\
好み追従率（1ビット） & 好みと方角が食い違うエピソードで $u_t$ が好みに従う割合 & 好み混合／段階II & 事前登録（経験80\%），結果を見た後の確認（経験50\%） \\
\midrule
\multicolumn{4}{@{}l}{\textbf{照合}} \\
正解率 BA & 2つの候補のうち予測誤差の和が小さい方を選んだ正解率（配置と好みの4つの組み合わせの平均） & 好み混合／自己視界：L1，他者特徴：MSE & 事前登録した主（自己視界基準），副（他者特徴基準） \\
好み追従率 & 好みと方角が食い違う88エピソードで選択が好みに従う割合 & 好み混合／他者特徴：MSE & 事前登録（経験50\%の結果の後，経験80\%の結果の前） \\
方角だけの基準 & A-2 の $t=1$〜50 の変位の方角で選んだときの BA（データだけで決まる） & 好み混合 & 事前登録外の記述 \\
\midrule
\multicolumn{4}{@{}l}{\textbf{全体}} \\
組をまたいだ統計 & 各組の最終5点の平均を1つの値とし，組を対にした差の平均と $t$ 分布の95\%信頼区間 & --- & 事前登録外の記述（符号無作為条件の6組は判定として登録） \\
\midrule
\multicolumn{4}{@{}p{0.97\linewidth}@{}}{評価データ：共通＝共通テスト（経験なし条件の評価用データ），緑好み・赤好み・好み混合＝それぞれの評価用データ．チェックポイント：MSE・L1＝段階Iの MSE 期間・L1 期間の最終5点，L1 の最後＝段階Iの L1 期間の最後のチェックポイント（エポック 200），段階II＝段階IIの最終5点．} \\
\bottomrule
\end{tabular}
